\documentclass{article}
\usepackage[T1]{fontenc}
\usepackage{hyperref}
\usepackage[polish]{babel}
\usepackage[utf8]{inputenc}
\usepackage{authblk}
\usepackage{indentfirst}
\usepackage[export]{adjustbox}
\usepackage{graphics}
\usepackage{array}
\usepackage{graphicx}
\graphicspath{ {./images/} }
\usepackage{wrapfig}
\usepackage[protrusion=false,letterspace=200]{microtype}


\title{DiscreteCalculator}
\author{}
\date{}

\begin{document}

\maketitle

\section{Szczegółowy opis projektu}
    \par{Głównym celem projektu jest stworzenie interaktywnej aplikacji z interfejsem graficznym służącej do badań i wizualizacji pewnych struktur dyskretnych. Narzędzie ma pozwalać użytkownikowi na budowanie tych obiektów (m. in. grafów) oraz obserwowanie efektów działania na nich pewnych algorytmów lub przekształceń.}  
    \\
    \par{Aplikacja zostanie zaprojektowana według architektury Model-View-Controller, co zapewni oddzielenie warstwy matematyczno-algorytmicznej od widoku użytkownika. 
    Nasz projekt z założenia będzie umożliwiał łatwe dodawanie obsługi kolejnych algorytmów bez modyfikacji rdzenia aplikacji, co stanowi pole do wykorzystania mechanizmu refleksji. 
    Będziemy szeroko korzystać z programowania wielowątkowego, aby operacje użytkownika w interfejsie graficznym oraz wykonywanie algorytmów mogły działać (z pewną ograniczoną swobodą) współbieżnie. 
    Gwoli utrzymania czystego kodu planujemy opierać się na interfejsach oraz hierarchiach dziedziczenia.}
    \\
    \par{Jedną z głównych funkcjonalności projektu jest narzędzie do rysowania (ważonych) grafów prostych lub skierowanych za pomocą myszy oraz przycisków (redo/undo, przyciski tworzące wybrane klasy grafów). Dla stworzonych w ten sposób grafów będzie można wywoływać algorytmy i ich wizualizacje. Przykłady obejmują: przeszukiwanie wgłąb, przeszukiwanie wszerz, znajdowanie najkrótszych ścieżek dla wybranych wierzchołków, znajdowanie maksymalnego przepływu w sieci (wraz z obsługą wyjątków), znajdowanie kolorowań o pewnych własnościach.}
    \\
    \par{Opcjonalne funkcjonalności obejmują możliwość badania i wizualizacji innych klasycznych problemów kombinatorycznych. Jedną z potencjalnych ścieżek rozwoju projektu jest wizualizacja obiektów związanych z liczbami Catalana oraz wskazywanie obiektów im odpowiadających przez (naturalne) bijekcje. Drugą opcjonalną funkcjonalnością może być wizualizacja podziałów liczb o zadanych własnościach przez diagramy Ferrersa wraz z możliwością odbijania ich względem głównej przekątnej. Kolejną możliwą funkcjonalnością jest wizualizacja grafów ciągów binarnych/podzbiorów zbiorów przez kody Graya z uwzględnieniem elastyczności oferowanej przez Generics (zbiory mogłyby zawierać dowolne elementy).}

\end{document}
